\documentclass[
	12pt, %Schriftgröße
	a4paper,
	bibliography=totoc, %Inhaltsverzeichniseinträge f+r Quellen
	numbers=noenddot, %Entfernt Punkt hinter Gliederungsnummern
	ngerman, %Sprachpaket
	headsepline, %Headertrennlinie
	%footsepline, %Footertrennlinie
	oneside %einseitiges Druckformat %%% Unterdrücken der leeren Seite nach Titelblatt
	]{scrbook} %Dokumentenklasse (Koma-Script)

\author{Justin Becker}
\usepackage{microtype}
\usepackage[T1]{fontenc}
\usepackage{float}
\usepackage[utf8]{inputenc}
\usepackage[ngerman]{babel}
\usepackage{url}
\usepackage{graphicx} %Bilder einfügen
%\usepackage{hyphenat}
%\usepackage{pdfpages} %PDF einfügen
\usepackage[a4paper, margin=1in]{geometry}
\usepackage[right]{eurosym} %Euro-Zeichen
\usepackage{amssymb}
\usepackage{setspace} % Zeilenabstand
\usepackage[ 
   colorlinks,        % Links ohne Umrandungen in zu wählender Farbe 
   linkcolor=black,   % Farbe interner Verweise 
   filecolor=black,   % Farbe externer Verweise 
   citecolor=black,   % Farbe von Zitaten 
   urlcolor=blue	  % Farbe von Links
   ]{hyperref} %Verlinkungen
\usepackage[figure]{hypcap}
\usepackage[ngerman]{translator}
\usepackage{blindtext} % Lorem-Ipsum-Plugin
\usepackage[acronym, nonumberlist]{glossaries} %% use after hyperref %Glossar-Paket laden
%\usepackage[
	%nonumberlist, %keine Seitenzahlen anzeigen
	%acronym,      %ein Abkürzungsverzeichnis erstellen
	%toc,          %Einträge im Inhaltsverzeichnis
	%section      %im Inhaltsverzeichnis auf section-Ebene erscheinen
	%]
%{glossaries}

\usepackage{csquotes}
\usepackage[]{bookmark}
\usepackage{scrhack}

\usepackage{listings,xcolor} %Codeanzeige
\usepackage[normalem]{ulem}
\useunder{\uline}{\ul}{}

\usepackage{chngcntr}
\counterwithout{figure}{chapter}
\counterwithout{table}{chapter}

\usepackage{subfiles}

\usepackage{biblatex}
\addbibresource{lit.bib}

\definecolor{dkgreen}{rgb}{0,.6,0}
\definecolor{dkblue}{rgb}{0,0,.6}
\definecolor{dkyellow}{cmyk}{0,0,.8,.3}

\lstset{
    numbers=left, 
    numberstyle=\tiny, 
    numbersep=5pt,
    breaklines=true,
    frame=single,
    language=sh,
    basicstyle=\ttfamily\fontsize{10}{12}\selectfont,
    keywordstyle    = \color{dkblue},
    stringstyle     = \color{red},
    identifierstyle = \color{dkgreen},
    commentstyle    = \color{gray},
    emph            =[1]{php},
    emphstyle       =[1]\color{black},
    emph            =[2]{if,and,or,else},
    emphstyle       =[2]\color{dkyellow}
} 

%%%%%%%%%%%%%%%%%%%%%%%%%%%%%%%%%%%%%%%%%%%%%%%%%%%%%
%%%%%%%%%%% Sonderformatierung
%%%%%%%%%%%%%%%%%%%%%%%%%%%%%%%%%%%%%%%%%%%%%%%%%%%%%

% Seitenabstände definieren
\geometry{verbose,tmargin=3cm,bmargin=2cm,lmargin=3cm,rmargin=3cm} 

% Hurenkinder und Schusterjungen verhindern (Ja, das heißt wirklich so!!!)
\clubpenalty= 10000
\widowpenalty= 10000
\displaywidowpenalty= 10000 

%% Bei Referenzen im Text wird jetzt bei allen Ebenen "Kapitel" vorgestellt, z.b. Kapitel 2, Kapitel 2.2, Kapitel 6.3.2
\addto\extrasngerman{
    \def\sectionautorefname{Kapitel}%
    \def\subsectionautorefname{Kapitel}%
    \def\subsubsectionautorefname{Kapitel}%
    }

% Vertikaler Abstand zwischen Ende Textblock - Ende Fußzeile --> Abstand der Seitenzahl von Rand erhöhen 
\setlength{\footskip}{10mm}

% Abstand vor/nach Überschriften verändern

\RedeclareSectionCommand[%
    beforeskip=0.5\baselineskip,
    afterskip=0.5\baselineskip,
]{chapter}

\RedeclareSectionCommand[%
    beforeskip=0.5\baselineskip,
    afterskip=0.5\baselineskip,
]{section}

\RedeclareSectionCommand[%
    beforeskip=0.1\baselineskip,
    afterskip=0.1\baselineskip,
]{subsection}

\RedeclareSectionCommand[%
    beforeskip=0.01\baselineskip,
    %%afterskip=0.2\baselineskip
]{paragraph}

\setlength{\abovecaptionskip}{4pt}  % 1pc=12pt 
\setlength{\belowcaptionskip}{0pt}
%\setlength{\textfloatsep}{4pt}
\setlength{\intextsep}{1pc}

%% Verkleinerung der Textgröße unter Abbildungen
\addtokomafont{caption}{\small}



% Den Punkt am Ende der Glossareinträge deaktivieren
\renewcommand*{\glspostdescription}{}

%Glossar-Befehle anschalten
%\makeglossaries{}
\makenoidxglossaries{}
\loadglsentries{glossar.tex}

% sorgt dafür, dass bei Leerzeile die Einrückung verhindert und stattdessen eine Leerzeile eingefügt wird % erspart bigskips und erhöht die Lesbarkeit im LaTeX-Text 
\KOMAoptions{parskip=full*}

% ändert Titelschriftart in Serifen-Normalschriftart
\addtokomafont{disposition}{\rmfamily} 




%%%%%%%%%%%%%%%%%%%%%%%%%%%%%%%%%%%%%%%%%%%%%%%%%%%%%
\newcommand{\studentName}{Justin Becker}
\newcommand{\matrnr}{1233455}
\newcommand{\type}{Praxistransferbericht}
\newcommand{\topic}{Einbindung von Code Compliance Checks in der SAP Software Entwicklung}
\newcommand{\studiengangh}{Informatik}
\newcommand{\fachbereich}{FB2: Duales Studium --- Technik}
\newcommand{\studiengang}{Informatik}
\newcommand{\company}{SAP SE}
\newcommand{\betreuerHS}{Lara Maria Stricker}
\newcommand{\betreuerUnt}{Jenny Peter}
\newcommand{\jahrgang}{2023}
%%%%%%%%%%%%%%%%%%%%%%%%%%%%%%%%%%%%%%%%%%%%%%%%%%%%%>>>>>>>


\newcommand{\quickwordcount}[1]{%
  \immediate\write18{texcount -1 -sum -merge -q -relaxed #1.tex > #1-words.sum }%
  \input{#1-words.sum}%
}

%%%%%%%%%%%%%%%%%%%%%%%%%%%%%%%%%%%%%%%%%%%%%%%%%%%%%%%%%%%%%%

\begin{document}

% falsche Default-Silbentrennung überschreiben
%%falsche Default-Silbentrennung überschreiben
\hyphenation{Soft-ware-ent-wick-lung}

\hyphenpenalty=10000 \spaceskip=0.3em plus 4em minus 0.2em

%%%%%%%%%%%%%%%%%%%%%%%%%%%%%%%%%%%%%%%%%%%%%%%%%%%%%>>>>>>>
%%%%%%%%%%% Titelblatt

%% Anordnung und Aussehen von Titel und Untertitel

\subject{\type}

\title{
\normalfont\endgraf\rule{\textwidth}{.4pt}
\begingroup
	\centering
	\linespread{1.5}
	\huge\topic%
\endgroup
\endgraf\rule{\textwidth}{.4pt}
}
 
%%Eigentlich nicht besonders schön, aber Koma erlaubt die Anordnung eines weiteren Felder (hier: Fachbereich) nicht
\date{\normalsize vorgelegt am 8. August 9834\\ \textbullet{} \\ Fachbereich Duales Studium Wirtschaft / Technik \\
Hochschule für Wirtschaft und Recht Berlin}
%% \date muss leer angegeben werden, um die Default-Datumsanzeige zu unterdrücken

\publishers{
	\begin{tabular}{l l}
	\textbf{\normalsize{Name:}} & \normalsize{\studentName}  \tabularnewline%
    \textbf{\normalsize{Matrikelnummer:}} & \normalsize{\matrnr}  \tabularnewline%
	\textbf{\normalsize{Ausbildungsbetrieb:}} & \normalsize{\company}  \tabularnewline%
    \textbf{\normalsize{Fachbereich:}} & \normalsize{\fachbereich} \tabularnewline%
    \textbf{\normalsize{Studienjahrgang:}} & \normalsize{\jahrgang} \tabularnewline%
	\textbf{\normalsize{Studiengang:}} & \normalsize{\studiengang}  \tabularnewline%
	\textbf{\normalsize{Betreuerin Unternehmen:}} & \normalsize{\betreuerUnt} \tabularnewline%
	\textbf{\normalsize{Betreuerin Hochschule:}} & \normalsize{\betreuerHS} \tabularnewline%
	\textbf{\normalsize{Wortanzahl:}} & \normalsize{\quickwordcount{main}}
	\end{tabular}
}

\titlehead{\begin{center}
    \includegraphics[scale=0.7]{bilder/header_logo.PNG}
    \end{center}
    }

\maketitle
\onehalfspacing% anderthalbfacher Zeilenabstand

%%%%%%%%%%%%%%%%%%%%%%%%%%%%%%%%%%%%%%%%%%%%%%%%%%%%%%%%%%%%%%%%%%%%%%%%%%%%%%%%%%%%%%%%%%%%%%%%%%%%%%%%%%%%%%%%%%%%%%%%%%%
%%%%%%%%%%% Dokumenteninhalt START
%%%%%%%%%%%%%%%%%%%%%%%%%%%%%%%%%%%%%%%%%%%%%%%%%%%%%%%%%%%%%%%%%%%%%%%%%%%%%%%%%%%%%%%%%%%%%%%%%%%%%%%%%%%%%%%%%%%%%%%%%%%

%%%%%%%%%%%%%%%%%%%%%%%%%%%%%%%%%%%%%%%%%%%%%%%%%%%%%
%%%%%%%%%%% Abstract
\chapter*{Abstract}
\addcontentsline{toc}{chapter}{Abstract}
\subfile{chapter/abstract.tex}

\pagenumbering{Roman} % römische Seitenzahlen

%%%%%%%%%%%%%%%%%%%%%%%%%%%%%%%%%%%%%%%%%%%%%%%%%%%%%
%%%%%%%%%%% Inhaltsverzeichnis, Tabellen, Abbildungen, etc.
\newpage

\tableofcontents{}
\addcontentsline{toc}{chapter}{Inhaltsverzeichnis}

\listoffigures
\addcontentsline{toc}{chapter}{Abbildungsverzeichnis}

\listoftables
\addcontentsline{toc}{chapter}{Tabellenverzeichnis}

\addcontentsline{toc}{chapter}{Akronyme}
\printnoidxglossaries{}

\clearpage

%% arabische Seitenzahlen
\pagenumbering{arabic}

%%%%%%%%%%%%%%%%%%%%%%%%%%%%%%%%%%%%%%%%%%%%%%%%%%%%%
%%%%%%%%%%% Einführung

\chapter{Einleitung}\label{ch:einleitung}
\subfile{chapter/einleitung.tex}

\chapter{Software Entwicklung bei SAP}\label{ch:sap_software_entwicklung}
\subfile{chapter/sap_software_entwicklung.tex}

\section{Cloud Produkte}\label{sec:cloud}
\subfile{section/cloud.tex}

\section{CI/CD}\label{sec:cicd}
\subfile{section/cicd.tex}

\chapter{Code Compliance}\label{ch:code_compliance}
\subfile{chapter/code_compliance.tex}

\section{Beispiel}\label{sec:compliance_beispiele}
\subfile{section/compliance_beispiele.tex}

\section{Checks}\label{se~:checks}
\subfile{section/checks.tex}

\subsection{Scans}\label{subsec:scans}
\subfile{subsection/scans.tex}

\subsection{Scan Ergebnis Bewertung}\label{subsec:scan_bewertung}
\subfile{subsection/scan_bewertung.tex}

\chapter{Herausforderungen}\label{ch:herausforderungen}
\subfile{chapter/herausforderungen.tex}

\section{Externe Tools}\label{sec:externe_tools}
\subfile{section/externe_tools.tex}

\section{Code Beispiel}\label{sec:code_beispiel}
\subfile{section/code_beispiel.tex}

\chapter{Fazit}\label{ch:fazit}
\subfile{chapter/fazit.tex}

%%%%%%%%%%%%%%%%%%%%%%%%%%
% Quellen
%%%%%%%%%%%%%%%%%%%%%%%%%

\printbibliography

%% \bibliographystyle{alpha} %% tu es nicht, niemals, das ist eklig, nicht einkommentieren

\chapter{Ehrenwörtliche Erklärung}
\addcontentsline{toc}{chapter}{Ehrenwörtliche Erklärung}
\subfile{chapter/ehrenwörtliche_erklärung.tex}

\end{document}

